% ch2-background.tex -- part of thesis_v3.tex; do not add \documentclass here.
% Draft with the thesis-chapter workflow. Unsupported claims must use \fillin{...}.

\chapter{Theoretical Background}
\label{ch:theory_background}

This chapter collects the material the rest of the thesis depends on. It first
describes the power delivery network (PDN) design problem in the specific form in
which this work attacks it, then the simulation flow that produces the training
data and the formal statement of the decoupling capacitor (decap) placement
problem. The second half covers the machine-learning background: variational
autoencoders and multi-modal fusion, physics-informed and surrogate modelling,
and gradient-based optimization in a learned latent space. The chapter closes
with related work and the research gap that motivates the framework.

Throughout the thesis a placement is written as a binary vector
$\mathbf{b} \in \{0,1\}^{52}$ over the board's decap slots, the budget is
$K = \lVert \mathbf{b} \rVert_0$, the PDN input impedance seen at the observation
port is $Z(f)$ sampled at $231$ frequency bins spanning $1$--$600$~MHz, and
$Z_{\mathrm{target}}(f)$ denotes the frequency-dependent target impedance mask
the design must stay below.

Numerical settings quoted in this chapter are those of the implemented
framework. Their source locations in the code base are recorded in the
implementation index appendix rather than in the text, so that the discussion
here stays with the concepts and the method.

\section{Power Delivery Networks and the Optimization Problem}
\label{sec:pdn_optimization}

\subsection{Impedance Behavior and Decoupling Strategies}
\label{subsec:impedance_decoupling}

A power delivery network has one job: hold the supply voltage at the load pins
inside a tolerance while the current drawn by that load changes. A switching
digital load does not draw a constant current. Each transition charges and
discharges internal capacitance, so the current at the power pins carries a large
transient component whose energy is spread over a wide frequency range rather
than concentrated at the clock rate. Write $I(f)$ for the spectrum of that
current and $Z(f)$ for the impedance the network presents at the port where the
load is attached. The voltage disturbance seen at that port follows the
frequency-domain form of Ohm's law,
\begin{equation}
\lvert V_{\mathrm{ripple}}(f) \rvert = \lvert Z(f) \rvert \cdot \lvert I(f) \rvert .
\label{eq:ohm_freq}
\end{equation}
Equation~\ref{eq:ohm_freq} is what makes impedance the design variable. The
requirement is stated on the ripple, but the current spectrum belongs to the
silicon and not to the board, so the only factor the board designer controls is
$Z(f)$.

The target-impedance methodology turns that observation into a specification.
An allowed ripple $\Delta V_{\max}$ and a worst-case transient current
$I_{\max}$ are fixed, and their ratio
\begin{equation}
Z_{\mathrm{target}} = \frac{\Delta V_{\max}}{I_{\max}}
\label{eq:target_impedance}
\end{equation}
becomes the impedance the network must stay below. The approach and the
associated capacitor-selection procedure for CMOS power distribution are due to
\cite{smith1999target}. Because the load has spectral content across a broad
band, a single scalar is not enough in practice and the bound is imposed as a
mask over frequency; \cite{bogatin2018signal} develops this as the
frequency-domain target impedance method (ch.~13.22). That is the form used
throughout this thesis: $Z_{\mathrm{target}}(f)$ over the $231$ bins spanning
$1$--$600$~MHz, with the mask itself described below.

\subsubsection*{The decoupling capacitor as a series RLC circuit}

At the frequencies of interest a decoupling capacitor is not well described by
its capacitance alone. The standard lumped model is a series RLC branch
combining the capacitance $C$, an equivalent series resistance $\mathrm{ESR}$
and an equivalent series inductance $\mathrm{ESL}$:
\begin{equation}
Z_{\mathrm{cap}}(\omega) = \mathrm{ESR} + j\omega\,\mathrm{ESL}
+ \frac{1}{j\omega C},
\qquad
\lvert Z_{\mathrm{cap}} \rvert
= \sqrt{\mathrm{ESR}^2 + \left(\omega\,\mathrm{ESL} - \frac{1}{\omega C}\right)^{2}},
\label{eq:cap_rlc}
\end{equation}
with $\omega = 2\pi f$. The magnitude has a minimum at the series self-resonant
frequency, where the inductive and capacitive reactances cancel:
\begin{equation}
f_{\mathrm{SRF}} = \frac{1}{2\pi\sqrt{\mathrm{ESL}\cdot C}} .
\label{eq:srf}
\end{equation}
At $f_{\mathrm{SRF}}$ the impedance of the branch equals its $\mathrm{ESR}$.
Below it the branch is capacitive and $\lvert Z_{\mathrm{cap}} \rvert$ falls with
frequency as $1/(\omega C)$; above it the branch is inductive and
$\lvert Z_{\mathrm{cap}} \rvert$ rises as $\omega\,\mathrm{ESL}$. The
consequence is that one capacitor value is only useful over a limited band around
its own self-resonance, which is why a real PDN uses several values and several
parts rather than one large capacitor. The multilayer ceramic capacitor model,
the equivalent series inductance and the loop-inductance approximation behind it
are treated in \cite{bogatin2018signal} (ch.~13.12--13.14).

\subsubsection*{Anti-resonance in parallel capacitor banks}

Connecting capacitors in parallel does not lower the impedance everywhere.
Consider two branches at a frequency at which one is already above its
self-resonance, and therefore inductive, while the other is still below its own,
and therefore capacitive. The parallel combination of an inductive and a
capacitive branch has a parallel resonance, and at that frequency the two
susceptances cancel in the opposite sense to Equation~\ref{eq:srf}: the
combination presents an impedance \emph{peak} rather than a dip, limited only by
the loss in the loop. The same mechanism acts between a capacitor bank and the
inductance of the package and of the plane path that feeds it. Anti-resonant
peaks, not the resonant dips, are what violate $Z_{\mathrm{target}}(f)$, and they
are the reason adding a capacitor can make a design worse at some frequencies
while improving it at others.

The mitigations are all ways of either moving a peak, damping it, or lowering the
inductance that creates it: spreading capacitor values so that the individual
self-resonances and the peaks between them are staggered rather than aligned;
choosing parts whose $\mathrm{ESR}$ is large enough to damp the parallel
resonance instead of leaving it sharp; adding more capacitors, which
\cite{bogatin2018signal} treats explicitly as a way of engineering a reduced
parallel-resonant peak (ch.~13.16--13.17); and placing the capacitors so that the
inductance between them and the load is small, which is taken up next.
\fillin{Pin the value-spreading and ESR-damping strategies to a specific passage
in \cite{bogatin2018signal} or replace them. Ch.~13.16--13.17 is confirmed from
the publisher's table of contents to cover combining capacitors in parallel and
reducing the parallel-resonant peak by adding more capacitors; the other two
strategies are stated here from general design practice and were not read from a
verified page.}

\subsubsection*{Mounting inductance and the spatial nature of the problem}

The $\mathrm{ESL}$ quoted for a part is only one contribution to the inductance
of the current loop it forms with the load. The capacitor pads, the mounting
vias down to the power and ground layers, and the path through the planes back to
the point being decoupled all add to that loop. Above self-resonance the branch
impedance is set by the total loop inductance $L_{\mathrm{loop}}$ rather than by
$\mathrm{ESL}$ alone, so $\lvert Z \rvert \approx \omega L_{\mathrm{loop}}$ and
the penalty for a large loop grows linearly with frequency. Two identical
capacitors mounted differently, or sitting at different distances from the IC
power pin, therefore do not decouple equally well. Optimizing the mounting of a
capacitor is the subject of \cite{bogatin2018signal} ch.~13.15, with the
loop-inductance approximation in ch.~13.14.

This is the physical reason the placement problem is spatial rather than a
counting problem. The vector $\mathbf{b} \in \{0,1\}^{52}$ selects both how many
capacitors are fitted and where they sit relative to the observation port, and
two placements with the same budget $K$ do not in general produce the same
$Z(f)$. It is also why the occupancy modality is encoded on a graph whose edges
follow board adjacency rather than being treated as an unordered set of $52$
bits (Section~\ref{subsec:vae_multimodal}).

At the upper end of the band the plane pair the capacitors connect to stops
behaving as a single lumped node and becomes a distributed structure with
resonances of its own, which also enter $Z(f)$
\cite{novak2007frequency}.
\fillin{Confirm against a physical copy of Novak and Miller that the book covers
plane-pair cavity resonance, and cite the specific chapter. The table of contents
could not be retrieved for this draft, so no statement about cavity modes, their
frequencies or their spatial pattern is made here, and this sentence should
either be substantiated or removed.}

\subsubsection*{Anti-resonance as a training prior}

The framework does not carry a circuit model of anti-resonance. What it carries
is a shape prior on the predicted spectrum: every impedance peak must be
preceded by a resonance dip. The penalty is evaluated on the raw log-impedance
channel as the mean of
$\text{peak\_score} \cdot \exp(-10 \cdot \text{prior\_dip})$. From bin
index~$190$ upwards the spectrum is monotonically increasing by construction, so
the term acts mainly in the low-frequency region. It is therefore a training
regularizer with a physical motivation, and this chapter presents it as such
rather than as a derived circuit relation.

Two design levers follow from the physics above. The first is which slots carry a
capacitor, which is the discrete variable this thesis optimizes. The second is
which part is placed in a slot. \fillin{State explicitly that the decap type is
fixed in the main study, and give the capacitance, ESR, ESL and part number
actually used in the ECADSTAR simulations. No parameter table exists: the
occupancy modality is binary presence only, and although the exploratory fork
\texttt{exp060\_multitype\_occ} names capacitor types it does not record their
electrical parameters. Settled by the ECADSTAR component library or the PEB
project.}

\subsubsection*{The frequency-dependent target mask}

The impedance requirement enters the framework as a fixed mask over the $231$
frequency bins, running from $0.8~\Omega$ at the low-frequency end to
$49.99999999999999~\Omega$ at the high-frequency end. Only the mask values are
supplied; no analytic form is fitted to them.

\fillin{Derivation and justification of this mask: why $0.8~\Omega$ at the low
end rising to $50~\Omega$ at the high end, which load current and ripple budget
it corresponds to, and who specified it. Nothing in the repository records how
the array was produced -- this is a gap the author must close from the original
design brief.}

\figplace{Target impedance mask $Z_{\mathrm{target}}(f)$ over the 231 bins, on log--log axes.}

\subsection{Power Integrity Simulation and Multi-modal Design Data}
\label{subsec:pi_simulation_multimodal}

\subsubsection*{Simulation flow and dataset generation}

Training data is produced by an end-to-end flow in six steps: sample and merge
placement combination tables; build a PEB project in ECADSTAR script order;
simulate the impedance spectrum and the PI-distribution heatmaps on Windows
through a headless batch invocation of the field solver; append the result to the
raw multi-frequency dataset; normalize, by a robust per-MHz $\log(1+x)$ transform
with an optional $K \le 30$ filter; and then train, evaluate or active-learn.
Row $i$ of a combination table corresponds to the ECADSTAR object
PI-$(i{+}1)$ in output order.

The conventional manual design loop that the framework is meant to replace runs
as follows: place capacitors, simulate the spectrum, and if peaks violate the
mask, run a spatial PI distribution, move capacitors nearer the hotspots, and
repeat. This description is taken from the project documentation and not from a
measured study of designer behaviour.

\fillin{PI solver configuration: solver type, meshing, frequency-sweep setup and
convergence criteria. The documentation covers only the batch invocation and file
paths, so this must come from the ECADSTAR project settings or the tool
documentation.}

\fillin{Wall-clock cost and compute cost of one ECADSTAR PI simulation. This
number is the quantitative motivation for building a surrogate at all, and no
artifact in the repository records it. A single timed batch run would settle it.}

\fillin{Board description: layer stackup, layer count, dielectric, physical
dimensions, plane geometry and the definition of the observation port. The
active-learning configuration names only the port IC1\_Port1 and the powerbus
Power\_GND; no geometry document exists.}

\subsubsection*{The three design modalities}

The framework consumes one design as three coupled views.

\emph{Occupancy.} The board carries $52$ decap slots laid out on a $7 \times 8$
grid with four invalid cells at $(0,3)$, $(0,4)$, $(6,3)$ and $(6,4)$. The slots
are labelled C1--C52 in a fixed order, and index $i$ of the placement vector
$\mathbf{b}$ refers to the $i$-th label in that order. The budget of a sample is
recovered from the vector itself as $K = \sum_i b_i$.

\emph{Impedance spectrum.} The spectrum has exactly $231$ points, and the length
is enforced at ingestion: a raw file of any other length is rejected. Inside the
model the spectrum carries an explicit channel dimension, so a batch of shape
$(B,231)$ is expanded to $(B,1,231)$ and the impedance decoder returns
$(B,1,231)$.

The frequency grid itself is a stored array of $231$ values running from
$1000000.0$~Hz through $1100000.0$~Hz to $600000000.0$~Hz. The spacing is not
uniform in $\log_{10}$: the first $\log_{10}$ step is $0.0414$ and the last is
$0.0073$.
\fillin{Record who defined the 231-bin sweep and on what basis, and state whether
it matches the ECADSTAR sweep settings exactly. No specification document defines
the grid; it exists only as the stored array.}

\emph{PI-distribution heatmap.} The spatial view is interpolated onto a
$64 \times 64$ grid. At creation it is a two-channel stack of the masked
impedance field and the board mask itself. The normalized arrays that reach
training are single-channel and of shape $(1,64,64)$, which is also the input and
output width of the heatmap encoder and decoder.

\figplace{The three modalities for one design: 52-slot occupancy grid, 231-bin spectrum, 64x64 PI heatmap.}

\subsubsection*{Frequency anchors and conditioning}

The spatial heatmap is frequency-dependent, so the dataset is built at a discrete
set of inspection frequencies called anchors. The current normalized training set
declares $24$ PI frequency anchors in MHz: $10.0$, $63.0$, $80.0$, $100.0$,
$120.0$, $150.0$, $170.0$, $180.0$, $200.0$, $230.0$, $250.0$, $270.0$, $280.0$,
$300.0$, $330.0$, $350.0$, $370.0$, $390.0$, $400.0$, $420.0$, $430.0$, $450.0$,
$470.0$ and $500.0$. The same list appears in the normalization statistics of
that dataset.

The dataset metadata is the authoritative anchor definition and is resolved
first at runtime. Two legacy definitions in the code base disagree with it -- a
built-in default tuple $(10, 63, 80, 130, 150, 200, 250, 270, 330, 400, 500)$
and a configuration file listing $20$ entries -- and both are reached only as
fall-backs, so a run that silently fell through to one of them would train
against a different anchor set.

Frequency enters the model as a scalar conditioning variable, normalized in
$\log_{10}$ between $1 \times 10^{6}$~Hz and $600 \times 10^{6}$~Hz and clipped
to $[0,1]$:
\begin{equation}
\mathrm{PI\_freq}(f) =
\mathrm{clip}\!\left(
\frac{\log_{10} f - \log_{10}(1{\times}10^{6})}
     {\log_{10}(600{\times}10^{6}) - \log_{10}(1{\times}10^{6})},\, 0,\, 1
\right).
\label{eq:pifreq_norm}
\end{equation}
For the current dataset the normalization constants are a $\log_{10}$ minimum of
$6.0$ and a $\log_{10}$ range of $2.778151250383644$. When a batch carries no
frequency, the normalized value for $200.0$~MHz is substituted.

\subsubsection*{Dataset scale and normalization}

The normalized training dataset contains $582997$ manifest rows over $24979$
unique layouts, held as a layout store, and was built with a maximum-budget
filter that drops layouts with $K > 30$ before statistics and outputs are
computed. Per-anchor sample counts are $24979$ for every anchor except
$80$~MHz, which has $8480$.

Heatmap normalization for this dataset is a robust per-MHz $\log(1+x)$ transform
in its unbounded variant, using per-MHz median and interquartile-range
statistics, a background value of $-3.1792$, and recorded clip bounds of
$-1.6791870594024658$ and $6.275787830352783$. The same three values appear in
the training configuration of \texttt{exp059\_capacity\_freq}; the unbounded
loader path applies no clip. Impedance normalization is a logarithmic
$z$-score with $\mu_{\log} = -1.039238452911377$ and
$\sigma_{\log} = 1.7831873893737793$, fitted over $24979$ layouts.

Training uses a layout-level holdout: all frequencies of one design identifier go
to train or to validation but never both, with a train split of $0.9$ and seed
$42$. \fillin{State why layout-level rather than row-level splitting matters
here in one sentence, with the leakage argument written out.}

\subsection{Formulation of the Decoupling Capacitor Placement Problem}
\label{subsec:decap_placement_formulation}

The engineering problem is to keep $Z(f)$ below a frequency-dependent target mask
over $1$--$600$~MHz sampled at $231$ bins, with the design lever being the binary
placement $\mathbf{b} \in \{0,1\}^{52}$ and $K = \lVert \mathbf{b} \rVert_0$.
Exhaustive search over the $2^{52}$ placements is intractable.

\subsubsection*{Feasibility with a safety margin}

The framework does not test feasibility against the mask directly. It tests
against the mask shifted inward by a constant safety margin. Writing
$\hat{Z}_j$ for the predicted impedance at bin $j$ and $Z_{\mathrm{target},j}$
for the mask value at that bin, both expressed in the comparison space, the
search counts a candidate feasible if and only if
\begin{equation}
\hat{Z}_j \;\le\; Z_{\mathrm{target},j} - \delta
\qquad \text{for all } j = 1,\dots,231,
\label{eq:feasibility}
\end{equation}
with a margin of $\delta = 0.1$.

Equation~\ref{eq:feasibility} is a surrogate-predicted feasibility test, not a
verified one. The quantity $\hat{Z}_j$ is the model's prediction, so a candidate
the search flags as feasible has not been simulated at that point: ground-truth
confirmation requires exporting the placement and re-running the ECADSTAR PI
simulation. Wherever this thesis calls a candidate feasible without saying
otherwise, the word refers to Equation~\ref{eq:feasibility} on predicted
impedance and is not a sign-off.

The comparison space matters for interpreting $\delta$. The search runs by
default in logarithmic space, in which $\hat{Z}$ and $Z_{\mathrm{target}}$ are
the $z$-scored log-impedance, so $\delta = 0.1$ is in normalized units and is not
$0.1~\Omega$. In that mode the mask is converted by $\log Z_{\mathrm{target}}$
followed by $z$-scoring with the dataset impedance log mean and standard
deviation.

\fillin{How the value $0.1$ was chosen, and what it corresponds to in ohms at a
representative frequency. Nothing in the repository records the rationale, and
its physical size depends on the $z$-scoring constants above.}

\subsubsection*{The per-budget objective}

The search is run per budget $K$. Let $Z^{\mathrm{bnd}}_j = Z_{\mathrm{target},j} - \delta$
be the shifted boundary, and define the per-bin violation
\begin{equation}
v_j = \big[\max(0,\; \hat{Z}_j - Z^{\mathrm{bnd}}_j)\big]^{p},
\qquad p = 2.0 .
\label{eq:violation}
\end{equation}
The violation is aggregated as the sum of its mean and its maximum over bins,
and a reward is granted for headroom below the boundary:
\begin{equation}
\mathrm{exceed} = \frac{1}{231}\sum_j v_j + \max_j v_j,
\qquad
\mathrm{gap} = \frac{1}{231}\sum_j \mathrm{below}_j ,
\label{eq:exceed_gap}
\end{equation}
where $\mathrm{below}_j$ is the headroom of bin $j$ below the boundary. The base
objective in the default mode is
\begin{equation}
\mathcal{L}_{\mathrm{base}}
= -\,w_{\mathrm{gap}} \cdot \mathrm{gap} \;+\; w_{\mathrm{exc}} \cdot \mathrm{exceed},
\qquad
w_{\mathrm{gap}} = 1.0, \quad w_{\mathrm{exc}} = 25.0 .
\label{eq:base_objective}
\end{equation}
The default objective mode is the gap-and-maximum form written above; in the
alternative feasibility-mapping mode the base term reduces to
$w_{\mathrm{exc}} \cdot \mathrm{exceed}$ only.

Among the candidates that satisfy Equation~\ref{eq:feasibility} on predicted
impedance, the default selection metric is the maximum predicted impedance in
ohms: the candidate with the lowest maximum is kept, and the best candidate is
only ever updated for seeds flagged feasible. If no seed is feasible for a
budget, the run records explicitly that no solution was found for that $K$ rather
than returning the least-bad infeasible design.

\subsubsection*{Scope of the implemented formulation}

The formulation above is the one implemented in the Stage-2 latent-optimization
script, which loads the VAE and the impedance surrogate of
\texttt{exp038\_true\_multi} for both roles, and whose normalization-statistics
path refers to a dataset directory that is not present in the current checkout
-- the live dataset is the unbounded normalized multi-frequency set described in
Section~\ref{subsec:pi_simulation_multimodal}. There is no Stage-2 entry point
wired to \texttt{exp059\_capacity\_freq}. The formulation is therefore presented
here as implemented, and no inverse-design result may be attributed to the
current surrogate on this basis.
\fillin{Either port Stage 2 to \texttt{exp059\_capacity\_freq} and update the
normalization-statistics path, or state in the thesis that Stage 2 is
demonstrated on \texttt{exp038\_true\_multi} only. Also write the per-$K$
objective against the exp059 model if that port happens.}

\section{Generative Modeling and Latent-Space Optimization}
\label{sec:generative_latent_opt}

\subsection{Variational Autoencoders and Multi-modal Learning}
\label{subsec:vae_multimodal}

\subsubsection*{The variational autoencoder}

The variational autoencoder \cite{kingma2013auto} learns a latent variable model
by jointly training an encoder $q_\phi(\mathbf{z}\mid\mathbf{x})$ and a decoder
$p_\theta(\mathbf{x}\mid\mathbf{z})$ against the evidence lower bound, trading
reconstruction quality against a Kullback--Leibler term that pulls the
approximate posterior toward a prior.

Written out, for an observation $\mathbf{x}$ and a latent vector
$\mathbf{z} \in \mathbb{R}^{d_{\mathrm{latent}}}$ with prior
$p(\mathbf{z}) = \mathcal{N}(\mathbf{0}, \mathbf{I})$, the marginal likelihood
$\log p_\theta(\mathbf{x})$ is intractable and training maximizes the evidence
lower bound (ELBO) instead:
\begin{equation}
\log p_\theta(\mathbf{x}) \;\ge\;
\mathcal{L}_{\mathrm{ELBO}}(\theta, \phi; \mathbf{x}) =
\mathbb{E}_{q_\phi(\mathbf{z}\mid\mathbf{x})}\!\left[\log p_\theta(\mathbf{x}\mid\mathbf{z})\right]
\;-\; D_{\mathrm{KL}}\!\left(q_\phi(\mathbf{z}\mid\mathbf{x}) \,\big\|\, p(\mathbf{z})\right),
\label{eq:elbo}
\end{equation}
where $\theta$ are the decoder parameters, $\phi$ the encoder parameters and
$D_{\mathrm{KL}}$ the Kullback--Leibler divergence. The first term is the
reconstruction term: it rewards latents from which the decoder can reproduce the
input. The second is a regularizer that pulls the approximate posterior
$q_\phi(\mathbf{z}\mid\mathbf{x})$ toward the prior. The gap between
$\log p_\theta(\mathbf{x})$ and the bound is
$D_{\mathrm{KL}}(q_\phi(\mathbf{z}\mid\mathbf{x}) \,\|\, p_\theta(\mathbf{z}\mid\mathbf{x}))$,
which is non-negative and vanishes only when the approximate posterior equals the
true one \cite{kingma2013auto}.

The encoder emits a mean $\boldsymbol{\mu}_\phi(\mathbf{x})$ and a diagonal
variance $\boldsymbol{\sigma}^2_\phi(\mathbf{x})$, and the sampling step is made
differentiable by the reparameterization trick: instead of sampling
$\mathbf{z}$ from $q_\phi$ directly, one draws
$\boldsymbol{\epsilon} \sim \mathcal{N}(\mathbf{0}, \mathbf{I})$ and sets
\begin{equation}
\mathbf{z} = \boldsymbol{\mu}_\phi(\mathbf{x})
+ \boldsymbol{\sigma}_\phi(\mathbf{x}) \odot \boldsymbol{\epsilon},
\label{eq:reparam}
\end{equation}
with $\odot$ the elementwise product, so that the gradient of the first term of
Equation~\ref{eq:elbo} with respect to $\phi$ passes through the sample
\cite{kingma2013auto}. Equations~\ref{eq:elbo} and~\ref{eq:reparam} are what
make the latent usable as a search domain later: the KL term keeps posteriors
from spreading arbitrarily far from the prior, so the region of
$\mathbb{R}^{d_{\mathrm{latent}}}$ that carries prior mass is the region the
decoder has been trained on. Section~\ref{subsec:latent_space_opt} relies on
this, and the Stage-2 objective carries an explicit prior term and a
posterior-boundary term for that reason.

The weighting of the KL term against reconstruction is the $\beta$-VAE
formulation \cite{higgins2017beta}, and conditioning a VAE on side information
follows the conditional generative framework of \cite{sohn2015conditional}.

In \texttt{exp059\_capacity\_freq} the KL term is configured with
\texttt{free\_bits} $=0.12$, $\beta$ annealing disabled at a final $\beta$ of
$0.03$, a per-expert KL weight of $0.02$, and a $\sigma$-regularization weight of
$1.5$ toward a target of $0.45$.

The first of those settings addresses a known failure mode of
Equation~\ref{eq:elbo}. The KL term is minimized outright by making
$q_\phi(\mathbf{z}\mid\mathbf{x})$ equal to the prior for every input. In that
regime the encoder carries no information about $\mathbf{x}$ and the decoder
ignores the latent altogether; the approximate posterior has collapsed onto the
prior, which is the failure mode named and analysed as posterior collapse in
\cite{he2019posterior}. For this framework the consequence is
specific rather than merely cosmetic: Stage 2 optimizes $\mathbf{z}$ against the
target mask, and a latent the decoder ignores cannot steer the decoded placement.

The countermeasure used here is free bits, or KL thresholding, introduced
alongside the inverse-autoregressive-flow work of \cite{kingma2016iaf}. The idea
is to stop penalizing the KL contribution of a latent dimension once it has
fallen below a fixed threshold, so that each dimension may retain a minimum
amount of information at no cost and the optimizer has no incentive to drive it to
zero. The threshold used here is the value $0.12$ quoted above.
\fillin{Confirm the exact hinge form of the free-bits objective and the units of
the threshold -- nats or bits, per latent dimension or per group of dimensions --
directly from the PDF of \cite{kingma2016iaf}, and check it against the exp059
implementation, before writing the equation. The formulation could not be read
from the paper's own full text for this draft, so only the qualitative
description above is given.}
\fillin{Explain what the $\sigma$-regularizer (weight $1.5$ toward a target of
$0.45$) is doing here, and why the free-bits threshold $0.12$, the final $\beta$
of $0.03$ and the per-expert KL weight $0.02$ were chosen at those particular
values; the configuration records the values but no note records the reasoning.}

\subsubsection*{Product-of-Experts fusion}

With several modalities describing one design, the joint posterior must be formed
from per-modality posteriors. The Product-of-Experts multimodal VAE
\cite{wu2018mvae} multiplies Gaussian experts, which for Gaussians is a
precision-weighted average; the mixture-of-experts alternative is developed in
\cite{shi2019mmvae}, and \cite{kutuzova2021defense} discusses the trade-off
between the two.

The implementation follows the PoE form directly. Given experts with means
$\mu_i$ and variances $\sigma_i^2$,
\begin{equation}
\sigma^{-2}_{\mathrm{poe}} = \sum_i \sigma_i^{-2},
\qquad
\mu_{\mathrm{poe}} = \sigma^{2}_{\mathrm{poe}} \sum_i \frac{\mu_i}{\sigma_i^{2}},
\label{eq:poe}
\end{equation}
and a unit-Gaussian prior expert with $\mu = 0$, $\log \sigma^2 = 0$ is always
appended as an extra expert; $\epsilon = 10^{-8}$ is added to each variance and
to the combined precision. Expert log-variances are clamped to $[-6.0, 1.0]$ at
each expert head and the fused log-variance is clamped to $[-4.0, 2.0]$ before
reparameterization.

\fillin{Recorded rationale for choosing PoE over MoE or another fusion scheme for
this problem. Nothing beyond a description of what PoE does was found in the
repository, so this is the author's own design argument to write.}

\subsubsection*{Structured latent and expert wiring}

The latent is partitioned rather than shared uniformly. Writing
$\mathbf{z} = [\mathbf{z}_{\mathrm{shared}} \mid \mathbf{z}_{\mathrm{peak}} \mid \mathbf{z}_{\mathrm{spatial}}]$,
the widths are
\begin{align}
d_{\mathrm{shared}} &= d_{\mathrm{latent}} - d_{\mathrm{hm\,priv}} - d_{\mathrm{imp\,peak}},
\label{eq:shared_dim}\\
d_{\mathrm{layout}} &= d_{\mathrm{shared}} + d_{\mathrm{imp\,peak}},
\label{eq:layout_dim}
\end{align}
where $d_{\mathrm{hm\,priv}}$ is the number of heatmap-private dimensions and
$d_{\mathrm{imp\,peak}}$ the number of impedance-peak dimensions. For
\texttt{exp059\_capacity\_freq} the total latent width is
$d_{\mathrm{latent}} = 128$, with $d_{\mathrm{hm\,priv}} = 40$ and a
conditioning width $d_{\mathrm{cond}} = 32$. The peak width is not set in that
configuration and takes its default of $d_{\mathrm{imp\,peak}} = 8$, alongside
default graph settings of hidden width $128$, three layers and dropout $0.1$.
Equations~\ref{eq:shared_dim}--\ref{eq:layout_dim} then give
$d_{\mathrm{shared}} = 80$ and $d_{\mathrm{layout}} = 88$.

The experts are wired to different slices of that latent. The heatmap expert
predicts the full latent vector; the occupancy expert predicts only the shared
dimensions; the impedance expert predicts the shared plus peak dimensions; and
dimensions an expert does not own are padded with the prior
($\mu = 0$, $\log \sigma^2 = 0$) before the product is taken. The heatmap encoder
input width is $64 + 2 d_{\mathrm{cond}}$ because it sees both $K$ and the
normalized frequency, while the occupancy and impedance heads take
$64 + d_{\mathrm{cond}}$ and see $K$ only. A dedicated frequency expert maps the
frequency embedding to the heatmap-private dimensions alone, with its shared and
layout dimensions padded by the prior; that expert is enabled by default and is
not disabled in the exp059 configuration.

\figplace{Latent partition and expert-to-latent wiring: which expert writes which dimensions.}

\subsubsection*{Graph encoders over slots and frequency bins}

Two of the three modalities are encoded with graph neural networks in the spirit
of \cite{kipf2016graph} and the message-passing formulation of
\cite{gilmer2017message}. The occupancy GNN operates on a fixed slot graph: the
$52$ slots are nodes, edges connect four-neighbours on the $7 \times 8$ grid,
self-loops are added, and the adjacency is symmetrically degree-normalized as
$D^{-1/2} A D^{-1/2}$. Node features are the occupancy value ($1$), a learned
slot embedding (half the hidden width) and the normalized row and column
position ($2$), with positions normalized by $6.0$ and $7.0$ respectively; the
graph is pooled by concatenating the mean and the max over nodes and projected to
an output width of $64$.

The spectrum is treated as a chain graph over its $231$ bins, with optional
self-loops and the same symmetric degree normalization. Its node features are the
bin value ($1$), a learned bin embedding (half the hidden width) and a normalized
bin position ($1$) taken from an evenly spaced grid on $[0,1]$, pooled the same
way to width $64$.

Each graph convolution layer computes a linear transform of the node features
plus a linear transform of the neighbour aggregate $A\mathbf{x}$, followed by
LayerNorm, a LeakyReLU with slope $0.1$ and optional dropout, implemented in
plain PyTorch with no external graph-learning dependency. Before the latent
heads, the two $64$-wide features are concatenated and projected from width $128$
to width $64$ with LayerNorm and LeakyReLU into a joint feature that feeds both
the occupancy and the impedance expert head.

Both discrete modalities are decoded by graph decoders as well: a global latent
and conditioning vector is broadcast to every node, concatenated with node
embeddings and positions, passed through GNN layers and read out per node --
$52$ logits for occupancy, $231$ values for the spectrum.

\subsubsection*{Conditioning on budget and frequency}

The budget is embedded with a learned table of $53$ entries, covering
$K \in \{0,\dots,52\}$, and the conditioning vector handed to the heatmap path is
the concatenation of the $K$ embedding and the frequency embedding, of width
$2 d_{\mathrm{cond}}$.

Frequency is embedded by a Fourier-feature MLP. The input is
$[f,\ \sin(2\pi f b),\ \cos(2\pi f b)]$ for $b = 1,\dots,B_{\mathrm{four}}$,
giving an input width of $1 + 2B_{\mathrm{four}}$, mapped by
Linear--SiLU--Linear--SiLU to $d_{\mathrm{cond}}$. In
\texttt{exp059\_capacity\_freq} the number of Fourier features is
$B_{\mathrm{four}} = 16$, where \texttt{exp057\_structured\_graph} and
\texttt{exp058\_asymmetric\_kl} use $8$.

Conditioning is injected into the heatmap decoder by feature-wise linear
modulation \cite{perez2018film}. The implementation applies
$x \mapsto x \cdot (1 + \gamma) + \beta$ with $\gamma$ and $\beta$ produced by a
Linear layer from the conditioning vector, whose weights and bias are
zero-initialized except that the bias entries for the $\gamma$ half are filled
with $1.0$.
\fillin{Reconcile the FiLM initialization: the forward pass adds $1.0$ to
$\gamma$ while the bias for the $\gamma$ half is also set to $1.0$, so the layer
does not start exactly at identity. Confirm the intended behaviour against the
exp043 frequency-conditioning module and state it.}

Multi-scale conditioning is enabled, so FiLM is instantiated at four decoder
scales: the $8\times8$ bottleneck on $128$ channels, $16\times16$ on $64$
channels, $32\times32$ on $32$ channels, and $64\times64$ on $16$ channels, each
conditioned on the $2 d_{\mathrm{cond}}$ vector; a further FiLM layer on $32$
channels sits in the heatmap encoder. The corresponding switch is not present in
the exp059 configuration file and defaults to enabled, whereas the
\texttt{exp057\_structured\_graph} model builder has no such key at all.

The exp059 heatmap decoder has no U-Net skip connections: it is driven by the
latent and by FiLM conditioning only, which is the form Stage 2 and unconditional
generation require. The \texttt{exp057\_structured\_graph} builder still exposes
a heatmap skip-connection switch, enabled by default; the exp059 builder has no
such key.

\subsubsection*{Modality dropout and the occupancy-only path}

At inference time, when a placement is being proposed, the spectrum and the
heatmap are exactly the quantities that are not yet known. The encoder must
therefore work from occupancy alone. Modality dropout is included so that the
encoder is trained on subsets of the modalities, which is the condition the
occupancy-only path faces at inference. During training each of the three data
experts is independently kept with probability $1 - p_{\mathrm{drop}}$ and
otherwise replaced by the prior; if all three are dropped, one is forced back on.
The frequency expert is appended after the dropout step and is never dropped, and
an optional protection flag forces the heatmap expert on. In
\texttt{exp059\_capacity\_freq} the dropout probability is
$p_{\mathrm{drop}} = 0.06$, both initially and as the start value of the
schedule, heatmap protection switches on at epoch $25$, corresponding to a
training fraction of $0.325$, and the focus phase uses $p_{\mathrm{drop}} =
0.03$.
\fillin{Quantify the effect of modality dropout on the occupancy-only path:
occupancy-only reconstruction and heatmap error with and without modality
dropout, from a matched pair of runs. No such ablation exists in the repository,
so this draft states the design intent behind the mechanism and makes no causal
claim that the path would be unusable without it.}

The occupancy-only encode is also available as an explicit inference path: the
impedance input is set to zero, the joint layout encoder is run, and the
posterior is formed from the occupancy expert together with the PoE prior expert
alone. This is the same path the active-learning loop uses for scoring and the
occupancy-only fine-tune uses for training.

During training the decode latent is sometimes taken from a layout encode instead
of the full posterior: with probability \texttt{layout\_train\_prob} the layout
branch is used, and within that branch with probability
\texttt{occ\_only\_encode\_prob} the occupancy-only encode is used. The full
encode is always run for the KL term; only the decode latent is swapped.

The live \texttt{exp059\_capacity\_freq} configuration sets
\texttt{layout\_train\_prob} $=0.65$ and \texttt{occ\_only\_encode\_prob} $=0.4$,
while the active-learning fine-tune override sets $0.9$ and $0.7$. The
experiment notes for the capacity study record both probabilities at $0$
(full-encode path only, active learning disabled) and both distillation weights
at $0$, whereas the current configuration sets the latent and output
distillation weights to $2.5$; the notes and the configuration also disagree on
the occupancy-only evaluation flag, the overlay data directory, the learning
rate, the number of epochs and the resume checkpoint.
\fillin{Resolve the disagreement between the exp059 notes and the live exp059
configuration by stating, per reported run, which configuration was actually in
force -- the checkpoint's own recorded config would settle it.}

The layout-to-private head is a two-hidden-layer MLP over the concatenation of
the $64$-wide occupancy feature, the $64$-wide impedance feature and the
$2 d_{\mathrm{cond}}$ conditioning vector. It predicts the heatmap-private latent
dimensions, optionally FiLM-modulated by the frequency embedding and squashed as
$4\tanh(x/4)$, and acts as the distillation student for the heatmap encode as
teacher. Both the private head and its frequency modulation are enabled by
default and neither is disabled in the exp059 configuration. The configured
distillation weights are $2.5$ for the latent term, a multiplier of $3.0$ on the
private dimensions, $0.3$ for the log-variance term and $2.5$ for the output
term -- subject to the disagreement recorded above.

\subsubsection*{Model capacity across the experiment lineage}

\texttt{exp057\_structured\_graph} and \texttt{exp058\_asymmetric\_kl} used a
latent width of $65$, a heatmap-private width of $15$, a conditioning width of
$8$ and $8$ Fourier features; \texttt{exp059\_capacity\_freq} raises these to
$128$, $40$, $32$ and $16$. exp058 differs from exp057 in the layout-path
training mix and the KL settings: \texttt{layout\_train\_prob} $0.9$ and
\texttt{occ\_only\_encode\_prob} $0.7$ against $0.65$ and $0.4$, a free-bits
threshold of $0.12$ against $0.08$, and a final $\beta$ of $0.03$ against $0.05$.
\texttt{exp060\_multitype\_occ} is a fork of exp059 that replaces binary
occupancy with a multi-type one-hot schema
$\mathrm{occ} \in \{0,1\}^{52 \times T}$, an all-zero row meaning empty, default
$T = 2$, decoded with sigmoid and focal binary cross-entropy rather than softmax
cross-entropy; that description comes from hand-written notes, not from a run
artifact, and exp060 carries the same four capacity values as exp059.

\subsection{Physics-Informed and Surrogate Modeling}
\label{subsec:physics_surrogate}

Surrogate modelling for signal and power integrity is an established direction;
\cite{swaminathan2020demystifying} and \cite{huang2021ml} survey the field for
packaging and for EDA more broadly. For PDN impedance specifically,
\cite{zhang2021fast} trains a deep network to predict impedance for arbitrary
board shape, stackup, IC location and decap placement, and
\cite{goay2022adaptive} models power--ground plane impedance with deep networks
under an adaptive sampling process. Worst-case power integrity has been predicted
with convolutional networks \cite{todaes2022worst}.

\fillin{One paragraph on what ``physics-informed'' means for this thesis
specifically: which constraints are imposed as loss terms versus which are
imposed structurally (hard top-$K$, board mask), and where that sits relative to
the physics-informed learning literature. This is an interpretive framing the
author should write.}

Three physics terms are implemented in \texttt{exp059\_capacity\_freq}. The
anti-resonance shape prior of Section~\ref{subsec:impedance_decoupling} is one.
The second and third come from a physics critic, a CNN that maps the
$7 \times 8$ decap occupancy grid together with the normalized inspection
frequency to a spatial effectiveness map in $[0,1]$ at $64 \times 64$. It is
supervised by inverting and min--max normalizing the heatmap foreground, and it
supplies a radius-influence penalty on reconstructed heatmap pixels that are cold
where decap coverage is weak. The critic is detached during the VAE step. The
configured weights are radius-influence $1.0$, critic supervision $2.0$ and
anti-resonance $0.5$.

Before the heatmap is decoded, occupancy is binarized to a hard top-$K$
zero/one vector, described in the source as CAD-aligned, optionally with a
straight-through estimator. The top-$K$ operation is vectorized through a
descending-sort threshold so that a per-row variable $K$ needs no Python loop.
\texttt{exp059\_capacity\_freq} enables the binary occupancy decode and disables
the straight-through variant during training.

Decoder input widths follow from the latent partition: the heatmap decoder
receives the full latent plus $2 d_{\mathrm{cond}}$ conditioning; the occupancy
decoder receives $\mathbf{z}_{\mathrm{shared}}$ plus the $K$ embedding; and the
impedance decoder receives $\mathbf{z}_{\mathrm{layout}}$ plus the $K$ embedding
plus a $64$-wide occupancy graph context, which is the passed occupancy if one is
given and otherwise the sigmoid of the reconstructed occupancy logits.

A frequency-weighted two-dimensional FFT loss on the heatmap is also implemented,
computing an $L_1$ distance between $\log(1+x)$ magnitudes of the real FFT of the
reconstruction and of the target, with an optional Gaussian bump in MHz. It is
not active in \texttt{exp059\_capacity\_freq}: its weight key is absent from the
configuration and from the configuration dataclass, so the lookup falls back to
$0.0$ and the term is skipped. The experiment notes state the opposite, namely a
weight of $1.5$, a mid-boost of $1.5$, a centre of $200$~MHz and a width of
$80$~MHz. The configuration is taken as authoritative, and this thesis does not
describe the spectral loss as an active training term.

\subsubsection*{Off-anchor evaluation}

Because heatmaps exist only at the $24$ anchors, generalization between anchors
has to be measured explicitly. \texttt{exp059\_capacity\_freq} declares
evaluation frequencies of $155.0$, $250.0$ and $265.0$~MHz with weights $2.5$,
$3.0$ and $2.5$ under the key \texttt{eval\_off\_anchor\_mhz}. That key name is
not literally accurate: $250.0$~MHz is itself one of the $24$ training anchors,
so only $155$ and $265$~MHz are genuinely off-anchor. Any evaluation reported
over the three-frequency set therefore mixes one in-distribution anchor with two
held-out frequencies, and Chapter~\ref{ch:results_evaluations} and
Chapter~\ref{ch:conclusion_outlook} must state which of the two readings is
meant.

Decoding at a non-anchor frequency is also available in blended form, by
log-blending the decodes at the two bracketing training anchors,
$\hat{h} = (1-w)\,\hat{h}_{\mathrm{lo}} + w\,\hat{h}_{\mathrm{hi}}$ with the
bracketing pair and the weight $w$ determined by the requested frequency. That
blend is a separate method; the training and evaluation path decodes directly at
the requested normalized frequency.

\subsubsection*{Active learning as a data-efficiency mechanism}

Because each label costs an ECADSTAR simulation, the framework retrains on
selected new layouts rather than on uniformly sampled ones. The general framing
is pool-based active learning \cite{settles2009active}; uncertainty for a deep
model can be obtained from MC dropout \cite{gal2016dropout} or from ensembles
\cite{lakshminarayanan2016deep}, and batch-mode selection for regression without
an ensemble is treated in \cite{batchal2022batch}. Gaussian-process bandit
optimization \cite{srinivas2010gp} is the reference point for the residual-GP
acquisition used in this work. The exp059 candidate generator searches budgets
$K$ from $3$ to $10$, with $1600$ candidates per $K$, $80$ worst per $K$, $640$
total ECAD simulations per cycle, a heatmap grid size of $64$, and an evaluation
frequency grid of $[120, 150, 155, 170, 180, 200, 230, 250, 265, 270, 280, 300]$
in MHz. The loop itself is described in
Chapter~\ref{ch:design_implementation}, not here.

\subsection{Optimization in Latent Representation Spaces}
\label{subsec:latent_space_opt}

The idea of optimizing a design by moving through a learned continuous latent
space and decoding the result was established outside electronics by
\cite{gomezbombarelli2018automatic}, which searches a VAE latent space of
molecules and decodes candidates back to discrete structures.

The argument for doing this is that it replaces a search over discrete objects
with a search over a continuous vector. The set $\{0,1\}^{52}$ of placements
supports no gradient: an objective evaluated on it can only be probed by flipping
bits according to some hand-chosen move set, and there is no direction in which
to step. Encoding designs into $\mathbb{R}^{d_{\mathrm{latent}}}$ and decoding
them back yields a domain in which a candidate can be moved by a differentiable
update instead. \cite{gomezbombarelli2018automatic} states this property
directly for molecules: continuous representations ``allow the use of powerful
gradient-based optimization to efficiently guide the search for optimized
functional compounds''. This thesis applies the same argument to decap placement,
subject to the qualification in the next paragraph.

The obstacle for PDN placement is that the decoded object is discrete. A
placement must be a $0/1$ vector with exactly $K$ ones before it can be handed to
CAD, and the top-$K$ selection that produces it is not differentiable. The
standard remedy is the straight-through estimator \cite{bengio2013straight};
continuous relaxations of discrete variables such as Gumbel-Softmax
\cite{jang2016gumbel} and the Concrete distribution \cite{maddison2016concrete}
are the alternative family.

The implementation uses a straight-through top-$K$ read-out. With $\mathbf{p}$
the soft occupancy probabilities and $\mathbf{b}_{\mathrm{hard}}$ the $K$-hot
vector obtained by selecting the $K$ largest entries,
\begin{equation}
\tilde{\mathbf{b}} = \mathbf{b}_{\mathrm{hard}} + \mathbf{p} - \mathrm{sg}(\mathbf{p}),
\label{eq:ste}
\end{equation}
where $\mathrm{sg}(\cdot)$ denotes the stop-gradient operator. The forward pass
therefore sees the hard $K$-hot vector while the backward pass passes identity
gradient to the probabilities. Without it the sorting operation behind the
top-$K$ selection severs the gradient and $\mathbf{z}$ never moves in response to
the spectrum target. This is the structural reason the estimator is required
here, and Section~\ref{subsec:research_gap} returns to it.

With the decoder frozen, the search runs Adam over $\mathbf{z}$ per budget $K$
with parallel seeds. The configured settings are budgets $K = 1,\dots,25$, $32$
candidate seeds, $1200$ steps, a learning rate of $5\times10^{-2}$ and gradient
clipping at $5.0$. The total objective adds to the base term of
Equation~\ref{eq:base_objective} a latent prior term, a spectrum shape
regularizer, a posterior-boundary term, a track term (MSE to target), the
anti-resonance physics term and a peak / dual-top-$K$ alignment term; with more
than one seed it further adds a diversity term and an occupancy-confidence term:
\begin{equation}
\begin{split}
\mathcal{L} = \overline{\mathcal{L}_{\mathrm{base}}}
&+ 0.05\,\mathcal{L}_{z\ell_2}
+ 2.0\,\mathcal{L}_{\mathrm{shape}}
+ 10.0\,\mathcal{L}_{\mathrm{post}}
+ 2.0\,\mathcal{L}_{\mathrm{track}} \\
&+ 0.5\,\mathcal{L}_{\mathrm{AR}}
+ w_{\mathrm{peak}}\,\mathcal{L}_{\mathrm{peak}}
\;(+\; 0.35\,\mathcal{L}_{\mathrm{div}}
+ 0.5\,\mathcal{L}_{\mathrm{conf}}),
\end{split}
\label{eq:stage2_total}
\end{equation}
with the posterior-boundary term limiting $\mathbf{z}$ to $2.5$ posterior
standard deviations.
\fillin{Insert the configured peak-loss weight $w_{\mathrm{peak}}$ from the
Stage-2 optimization script, and give the functional form of the shape, track,
posterior-boundary, peak, diversity and occupancy-confidence terms. Only the
weights were extracted, not the formulas.}

Two implementation details qualify what Equation~\ref{eq:stage2_total} optimizes
against. The impedance path defaults to a separate occupancy-to-spectrum
surrogate rather than the VAE impedance head; that surrogate is a direct map from
an occupancy vector of length $52$ to a log-$z$ spectrum of shape $(1,231)$ with
hidden widths $(256, 512, 512, 512, 256)$, and it does not use the normalized
inspection frequency at all. And, as recorded in
Section~\ref{subsec:decap_placement_formulation}, the whole script is wired to
\texttt{exp038\_true\_multi}.

\figplace{Stage-2 loop: z init, decode to soft occupancy, STE top-K, surrogate spectrum, loss against the mask, Adam step.}

\section{Related Work and Research Gap}
\label{sec:related_work}

\subsection{Power Delivery Network Optimization Methods}
\label{subsec:pdn_opt_methods}

\subsubsection*{Target impedance and heuristic search}

Target-impedance-based methodology and capacitor selection for CMOS PDNs is
described in \cite{smith1999target}. Board-level decap placement has been
attacked repeatedly with heuristic and evolutionary search. The titles record a
nature-inspired algorithm \cite{electronics2019nature}, a genetic algorithm
minimizing the number of capacitors \cite{ga2020minimum}, and an iterative
genetic algorithm combined with machine learning
\cite{electronics2020iterative}; simulation-driven module generation for decap
placement is reported in \cite{ars2023ai}. Only the bibliographic titles of those
four could be verified for this draft, so nothing is asserted here about the
evaluator each one queries.
\fillin{Confirm from each of electronics2019nature, ga2020minimum,
electronics2020iterative and ars2023ai what the evaluator is (a field solver or
a learned model) and whether it is queried for a scalar score without a gradient
being taken through it. Only then can a shared black-box-search structure be
claimed for this group.}

The one work in this group whose search mechanism could be verified is
\cite{convga2024gcn}, and it is also the closest complete instance to the present
framework: a graph-convolutional impedance surrogate wrapped in a genetic
algorithm that queries the surrogate's scalar value, with no gradient taken
through the surrogate. The distinction to hold onto is not the surrogate, which
both approaches have, but what is done with it -- a GA samples and scores,
whereas this thesis differentiates through the surrogate and moves a latent
vector.

\subsubsection*{Impedance surrogates}

\cite{zhang2021fast} predicts PDN impedance for arbitrary board shape, stackup,
IC location and decap placement, trained on more than one million
boundary-element-method boards. Its output is an impedance prediction; the
placement is an input to it, not something the model proposes.
\cite{goay2022adaptive} models power--ground plane impedance with deep networks
and an adaptive sampling process, and \cite{todaes2022worst} predicts worst-case
power integrity with a CNN. Those two are cited here from their titles only.
\fillin{Confirm from goay2022adaptive and todaes2022worst whether either
proposes a placement rather than only predicting a quantity, or drop the joint
reading of this group as prediction-only surrogates.}

\subsubsection*{Reinforcement learning for decap placement}

A distinct line, largely from KAIST, learns a placement \emph{policy} over
discrete actions. \cite{park2020drldecap} applies deep RL to decap design on
silicon interposers for 2.5D/3D ICs, meeting a target impedance with minimum
area. \cite{kim2021epeps} trains a pointer-network policy with REINFORCE as a
sequence-to-sequence model over candidate decap slots, handling arbitrary probing
ports and keep-out regions. \cite{park2022transformer} uses an attention-based RL
policy that parameterizes decap choice at multiple ports of an HBM PDN, and
reports minutes against more than a month for a full search.
\cite{kim2023devformer} generalizes the earlier pointer-network work of the same
first author into a symmetric transformer with permutation-symmetry and
relative-position inductive biases, and reports more than $30\%$ fewer decaps
than prior art on real hardware. Those two figures are as reported by the
respective papers; they were obtained on those authors' own boards and are not
comparable to any number in this thesis.

Outside that group, \cite{duan2024hierarchical} reports hierarchical decap
optimization for the PDN of 2.5D ICs based on deep reinforcement learning.
\fillin{Verify duan2024hierarchical against the paper itself. Only its
bibliographic title was available for this draft, so no statement is made here
about its action space, reward, evaluator or how the frequency- and time-domain
co-analysis enters the method.}

\subsection{Generative and Learning-Based Design Approaches}
\label{subsec:generative_learning_approaches}

Predicting a spatial field of the power network as an image is now a small
literature of its own. \cite{chhabria2021iredge} uses encoder--decoder networks
to produce on-chip static and dynamic IR-drop maps conditioned on layout,
which legitimizes decoding a PI field as an image rather than as a scalar.
\cite{zhao2024pdnnet} combines a GNN over the PDN with a CNN for dynamic IR-drop
prediction, which is independent external evidence that the graph structure of
the network carries information a purely spatial view does not.
\cite{gif2026conditional} is architecturally the closest neighbour of this
thesis's heatmap decoder: a conditional multimodal generative framework fusing
image and graph features into a conditional diffusion process for spatial IR-drop
maps. Two differences should be stated precisely rather than glossed: it is a
diffusion model where this thesis uses a VAE decoder, and it conditions on
layout where this thesis conditions on inspection frequency. All three operate on
on-chip IR drop, whereas the modality here is a board-level PI distribution at a
given frequency.

Learning-based EDA more broadly is surveyed in \cite{huang2021ml} and, for
packaging-level SI/PI, in \cite{swaminathan2020demystifying}. Placement as a
learned decision problem at chip level appears in \cite{mirhoseini2021chip}, and
graph/transformer models over interconnect in \cite{traceformer2024graph}.

\subsection{Research Gap and Motivation for the Proposed Framework}
\label{subsec:research_gap}

Across the works surveyed above, no paper was found that combines a multi-modal
generative surrogate over occupancy, spectrum and spatial PI distribution with
decoder-mediated latent-space gradient optimization through a discrete top-$K$
read-out for PDN decap placement. That statement is a claim about the literature
searched for this thesis, not a claim of superiority over any of it: none of the
comparisons below is a performance comparison, and no number in this thesis is
benchmarked against a number in those papers.

\subsubsection*{Distinction from the differentiable-surrogate line}

The methodologically closest work is the differentiable-surrogate line from the
examiner's own group. \cite{withoft2026amortized} performs amortized neural
optimization for pre-layout signal-integrity design-space exploration, mapping a
context to near-optimal design parameters in a single deterministic forward pass
through a differentiable surrogate. \cite{withoft2026buffer} builds
buffer-parameterized surrogate models over a $44$-parameter design space and
benchmarks neural, Gaussian-process and tree-based surrogates for cross-technology
SI analysis and optimization. \cite{ecik2026emd} takes a deliberately
non-gradient-based route: a deterministic decision-tree filter followed by an
Earth Mover's Distance ranking over continuous design-parameter predictions,
explicitly avoiding both inverse modelling and stochastic search.

The distinction rests on a targeted reading of all three, checking each for any
mention of a latent or embedding space, of discrete design variables, of top-$K$
or cardinality-constrained selection, and of a straight-through or relaxed
gradient estimator. None of the three contains any of them: the optimization
variable in each case is a vector of continuous parameters in the original
design-parameter space -- equalizer and buffer parameters, PCB parameters,
interconnect design variables. This thesis instead optimizes a latent vector
$\mathbf{z}$ that is decoded to a placement, and the read-out from the decoder to
the placement is a hard top-$K$ selection over $52$ discrete slots. That is
precisely why the straight-through estimator of Equation~\ref{eq:ste} is
structurally necessary here and is absent there: with a continuous parameter
vector there is nothing between the optimizer and the surrogate to break the
gradient, whereas here the top-$K$ selection would sever it.

Two further differences are worth stating without inflating them. Those works
address signal integrity; this one addresses power integrity, which the thesis
scope restricts it to. And the amortized formulation of
\cite{withoft2026amortized} produces a design in one forward pass, whereas the
Stage-2 formulation here runs an iterative search per budget $K$
(Section~\ref{subsec:latent_space_opt}).
\fillin{Say explicitly whether the amortized single-pass property is a
disadvantage of this thesis's approach in inference cost, and if so by how much.
That needs a timing measurement of the Stage-2 loop, which has not been recorded.}

\subsubsection*{Distinction from the reinforcement-learning paradigm}

The KAIST line -- \cite{park2020drldecap} on interposers, \cite{kim2021epeps}
with a pointer network and REINFORCE, \cite{park2022transformer} on HBM, and
\cite{kim2023devformer}, which generalizes the first of those two given the
shared first author -- is a third paradigm, distinct from both the
differentiable-surrogate line and from this thesis. Those methods learn a
placement policy over discrete actions or action sequences. There is no gradient
search through a differentiable surrogate, and there is no generative latent
space to search in at all: the model outputs actions, not a representation that
is decoded. The practical consequence is a difference in what has to be redone
when the requirement changes. For \cite{kim2023devformer} the literature survey
for this thesis records a policy that must be retrained per target scenario
class. On the other side, the framework here retargets to a new
$Z_{\mathrm{target}}(f)$ by re-running the latent search with the decoder
frozen, as implemented in the Stage-2 script against
\texttt{exp038\_true\_multi}. That retargeting has not been exercised against
the exp059 surrogate, for the wiring reason given in
Section~\ref{subsec:decap_placement_formulation}, so it is a property of the
implemented Stage-2 script and not a reported result.
\fillin{Confirm from park2020drldecap, kim2021epeps and park2022transformer
whether a change of target impedance mask requires policy retraining, and cite
the specific statement in each. Only kim2023devformer is recorded as requiring
per-scenario retraining, so the retargeting contrast is deliberately not
generalized to the whole reinforcement-learning line here.}

\subsubsection*{The closest use of latent representations in SI/PI}

The nearest hit in the SI/PI literature for the phrase ``latent
representations'' combined with optimization is \cite{usama2025dram}. It learns
latent representations of DRAM equalizer eye diagrams as a fast evaluation
feature, then runs an advantage actor--critic agent that optimizes the raw
equalizer parameters using that latent representation as a cheap fitness signal.
The latent space there is a frozen evaluation feature consumed by a policy over
raw parameters. It is never itself the optimization variable, and nothing is
decoded back through it. That is the exact distinction this thesis rests on, and
it holds even for the closest neighbour found.

\subsubsection*{Distinction from surrogate-plus-search and from spatial-map models}

Against \cite{convga2024gcn}, the difference is the use made of the surrogate:
a genetic algorithm queries it for a scalar and searches around it, whereas here
the gradient is taken through it. Against \cite{zhang2021fast}, the difference is
in the data regime: it generalizes across many boards using cheap boundary-element
labels, whereas this thesis trains one board deeply on more expensive ECADSTAR
labels, which is the direct motivation for the active-learning loop rather than
for broader sampling. Against the IR-drop spatial-map line
\cite{chhabria2021iredge, zhao2024pdnnet, gif2026conditional}, the difference is
both the quantity modelled -- board-level PI distribution at a chosen inspection
frequency, against on-chip IR drop -- and the role of the model: those works
predict a field, whereas here the spatial decoder is one branch of a joint
generative model whose latent is subsequently optimized.

\subsubsection*{Absence of a shared benchmark}

No shared benchmark exists across any of these PDN decap-placement works. The
Götze-group papers, the KAIST RL line, DevFormer, \cite{convga2024gcn} and
\cite{zhang2021fast} each use their own board, port definition and target mask.
None of their reported numbers is therefore directly comparable to any result in
this thesis, and no such comparison is made. A fair quantitative comparison would
require re-implementing at least one of those baselines on the $52$-slot board and
the $231$-bin mask used here.
\fillin{Decide and state which single baseline, if any, will be re-run on this
board for Chapter~5, and what the equal-budget protocol would be. As of this
draft no baseline has been re-run.}

\subsubsection*{From the gap to the research questions}

\fillin{The author's own statement of the research questions and hypotheses that
follow from this gap, agreed with the supervisor. This is deliberately not
drafted: the questions the thesis commits to answering are not recorded anywhere
in the repository.}

\fillin{One paragraph connecting the gap to the three contributions claimed in
Chapter~1, once those are finalized, with a cross-reference to the
corresponding chapter for each.}


%=============================================================================
% Chapter 3: Design and Implementation
%=============================================================================

% ---------------------------------------------------------------
% FILL IN checklist:
%   - Anti-resonance mitigations: value spreading and ESR damping are written from
%     general design practice. Pin them to a specific passage in bogatin2018signal or
%     replace them. Only ch.13.16-13.17 (combining capacitors in parallel; reduced
%     parallel-resonant peak by adding more capacitors) is confirmed from the TOC.
%   - Plane-pair cavity resonance: one sentence cites novak2007frequency, whose TOC
%     could NOT be fetched. Confirm from a physical copy and cite the chapter, or
%     delete the sentence. No cavity-mode behaviour is asserted in the draft.
%   - Decap type fixed in main study? give capacitance/ESR/ESL/part number used in
%     the ECADSTAR sims. Settled by: the ECADSTAR component library / PEB project.
%   - Derivation of the target mask (0.8 -> 50 ohm), stored as configs/target_impedance.npy.
%     Settled by: the original design brief; nothing in the repo records it.
%   - PI solver settings (mesh, sweep, solver type, convergence). Settled by: the
%     ECADSTAR project settings.
%   - Wall-clock cost of one ECADSTAR PI simulation. Settled by: one timed batch run.
%   - Board stackup/layers/dielectric/dimensions/port definition. Settled by: a
%     geometry document that does not yet exist.
%   - Provenance of the 231-bin sweep in configs/Frequency_data_hz.npy.
%   - Why layout-level (not row-level) holdout matters: leakage argument, one sentence.
%   - How the feasibility margin 0.1 (BOUNDARY_MARGIN) was chosen and its size in ohms.
%   - Port Stage 2 to exp059 (or state it is exp038-only) + fix NORMALIZATION_STATS_PATH.
%   - Free-bits formulation: exact hinge form and the units of the threshold (nats or
%     bits; per dimension or per group) from the kingma2016iaf PDF, cross-checked against
%     the exp059 implementation. NOT written as an equation in this draft on purpose.
%   - Meaning of the sigma-regularizer (weight 1.5 -> target 0.45) and why free_bits 0.12,
%     beta_final 0.03 and per_expert_kl_weight 0.02 were chosen at those values.
%   - Rationale for PoE over MoE for this problem. Settled by: author's design argument.
%   - FiLM init reconciliation: forward adds 1.0 to gamma AND bias gamma-half = 1.0.
%   - Resolve exp059 notes.md vs live config.yaml (layout_train_prob, occ_only_encode_prob,
%     distill weights, lr, epochs, resume). Settled by: the checkpoint's recorded config.
%   - "Physics-informed" framing paragraph: loss-imposed vs structurally-imposed constraints.
%   - PEAK_LOSS_WEIGHT value + functional forms of shape/track/posterior/peak/diversity/
%     occ-confidence terms. Settled by: pipelines/latent/optimize.py lines 330-416.
%   - Modality-dropout ablation: occupancy-only reconstruction/heatmap error with and
%     without modality dropout, from a matched pair of runs. No such run exists; until it
%     does, the chapter states design intent only and makes no causal claim.
%   - Evaluator/gradient status of electronics2019nature, ga2020minimum,
%     electronics2020iterative, ars2023ai. Settled by: reading the four papers. Until then
%     no shared black-box-search structure is claimed for that group.
%   - Whether goay2022adaptive or todaes2022worst proposes a placement (both cited from
%     title only). Settled by: reading the two papers.
%   - duan2024hierarchical method details (action space, reward, evaluator, how the
%     frequency/time-domain co-analysis enters). Only the title was verified.
%   - Inference-cost comparison vs the amortized single-pass formulation. Settled by:
%     a timing measurement of the Stage-2 loop (not recorded).
%   - Whether park2020drldecap, kim2021epeps and park2022transformer require retraining on
%     a new target mask; needs a quotation from each. Only kim2023devformer is recorded.
%   - Which baseline (if any) will be re-run on this board for Ch.5, and the protocol.
%   - Research questions/hypotheses following from the gap (author + supervisor).
%   - Paragraph linking the gap to the Ch.1 contributions.
%   - APPENDIX DEPENDENCY: the chapter preamble now promises an "implementation index
%     appendix". All file:line evidence removed from this chapter in the academic-register
%     pass is preserved in scratchpad thesis_phase1/ch2-impl-index.json (69 entries) and
%     must be rendered as that appendix, or the preamble sentence must be deleted.
%
% Citation verification queue (refs.bib note = "VERIFY: bibliographic fields not
% confirmed from source page") -- keys cited in this chapter that need checking:
%   smith1999target, settles2009active, maddison2016concrete, higgins2017beta,
%   electronics2019nature, shi2019mmvae, electronics2020iterative, ga2020minimum,
%   swaminathan2020demystifying, huang2021ml, mirhoseini2021chip, batchal2022batch,
%   goay2022adaptive, todaes2022worst, ars2023ai, convga2024gcn, traceformer2024graph
%   Additionally:
%   - bogatin2018signal (added for the PDN-physics gaps): carries a VERIFY note in
%     refs.bib. Publisher TOC confirms the ch.13 section numbers cited in prose
%     (13.12-13.17, 13.22), but two ISBNs appear across retailers for what looks like the
%     same 3rd edition, and the publication month is unconfirmed. Check the copy used and
%     fix the refs.bib fields before the first full build.
%   - novak2007frequency: carries a VERIFY note in refs.bib. Bibliographic data confirmed
%     via Google Books; the TOC could NOT be fetched (self-signed cert / 403), so the
%     book's coverage of plane-pair cavity resonance is UNVERIFIED. Cited once, with a
%     \fillin attached at the use site.
%   - kingma2016iaf: carries a VERIFY note in refs.bib. Title/authors/year confirmed from
%     the arXiv abstract page, but the free-bits definition could NOT be extracted from the
%     paper's full text. Cited for the idea only; no equation is written.
%   - smith1999target: the abstract is paywalled and was not read directly. Cited for the
%     target-impedance methodology, which is what it is universally cited for. No specific
%     numerical claim or derivation is attributed to it. The vault note Smith1999*.md links
%     a wrong DOI; the DOI in refs.bib (10.1109/6040.784476) is the correct one.
%   - he2019posterior: verified (ICLR 2019 / arXiv:1901.05534); no VERIFY note needed.
%   - convga2024gcn: literature block records verification as PARTIAL -- author list,
%     exact venue and page numbers unconfirmed (IEEE Xplore 10705608 paywalled).
%     Do not attribute authors in prose until confirmed.
%   - refs.bib AUTHOR-FIELD DEFECT (bibliography, not prose): convga2024gcn,
%     traceformer2024graph and ars2023ai carry placeholder strings in the author field
%     ("(IEEE conference paper 10705608)", "(DAC 2024)", "Advances in Radio Science,
%     vol. 21"). These will render as author names in the built PDF. Fix in refs.bib
%     before the first full build; the chapter deliberately attributes no authors in prose.
%   - Title-only citations. The following are cited from the title alone; the chapter says
%     so in the prose and asserts nothing about their mechanism: electronics2019nature,
%     ga2020minimum, electronics2020iterative, ars2023ai, goay2022adaptive,
%     todaes2022worst, duan2024hierarchical.
%   - kim2021epeps: title/authors/year/venue confirmed from the KAIST EE lab project
%     page, not from the IEEE Xplore landing page.
%   - park2020drldecap: confirmed via pure.kaist.ac.kr, not from IEEE Xplore directly.
%   - kim2021epeps and usama2025dram were added to refs.bib in an earlier drafting session.
%   - duan2024hierarchical, kutuzova2021defense, srinivas2010gp, gal2016dropout,
%     lakshminarayanan2016deep, bengio2013straight, jang2016gumbel, kipf2016graph,
%     gilmer2017message, perez2018film, wu2018mvae, kingma2013auto, sohn2015conditional,
%     gomezbombarelli2018automatic: cited from the "known" list of the literature block;
%     confirm each is present in refs.bib before the first full build.
%   - Hong/Li/Guan/Liang, "Impedance analysis and optimization research of power
%     distribution network" (Discover Applied Sciences 2025) was found but could NOT be
%     verified (paywall/403) and is deliberately NOT cited anywhere in this chapter.
%
% Evidence warnings:
%   1. docs/framework-overview.md sec.3 still calls exp057_structured_graph the current
%      surrogate stage; docs/experiment-lineage.md and CLAUDE.md say exp059_capacity_freq.
%      The academic-register pass REMOVED the paragraph that argued this out in the body
%      (it was repository meta). The document is treated as stale ONLY on which experiment
%      is current; its section-1 material is still used twice -- the conventional design
%      loop (subsec:pi_simulation_multimodal) and the engineering-problem statement incl.
%      2^52 (subsec:decap_placement_formulation) -- and the shape facts behind the second
%      (52 slots, 231 bins, frequency grid) are independently confirmed in source.
%   2. exp059 notes.md vs live config.yaml disagree on layout_train_prob (0 vs 0.65),
%      occ_only_encode_prob (0 vs 0.4), latent/output distill weights (0 vs 2.5),
%      eval_use_occ_only_layout, al_overlay_data_dir, learning_rate, num_epochs,
%      resume_checkpoint. Surfaced in one plain sentence, not resolved.
%   3. The frequency-weighted 2D-FFT spectral heatmap loss is IMPLEMENTED but its weight
%      key is absent from config.yaml and from the Config dataclass, so it falls back to
%      0.0 and the function returns None -- the term is DISABLED. exp059 notes.md claims
%      weight 1.5. This chapter does not describe it as an active loss term.
%   4. eval_off_anchor_mhz = [155, 250, 265] but 250 MHz IS one of the 24 training
%      anchors. Only 155 and 265 are genuinely off-anchor. Flagged in subsec:physics_surrogate.
%   5. src_vae/others/multifreq_anchors.py default tuple and configs/multifreq_anchors.yaml
%      (20 entries) both differ from the dataset's 24 anchors; code marks them legacy, so
%      dataset_meta.json wins. Flagged in subsec:pi_simulation_multimodal.
%   6. pipelines/latent/optimize.py is wired to exp038_true_multi and its
%      NORMALIZATION_STATS_PATH points at datasets/data_multifreq_norm/, which does not
%      exist. Stage 2 is presented as implemented, NEVER as demonstrated on exp059.
%   7. Provenance caveat: exp059 config.yaml, vae_multi_input_simple.py and
%      docs/framework-overview.md are untracked in git; the current exp059 architecture
%      and config are not committed, and the latest commit message still describes exp057
%      as current. Any line number in the implementation index may move.
%   8. Facts block lists 10 "missing" items covering board / solver / cost / part
%      parameters; every one is still marked with a \fillin rather than written from
%      background knowledge. The PDN-physics material is written from literature only
%      (bogatin2018signal, smith1999target) and is explicitly NOT sourced from the repo;
%      the anti-resonance term remains framed as a training regularizer with a physical
%      motivation, not as a derived circuit relation.
%   9. No experimental outcome is stated anywhere in this chapter, by design.
%  10. Equation eq:feasibility is a SURROGATE-PREDICTED test. The chapter states after it
%      that a candidate flagged feasible has not been simulated and that ground-truth
%      confirmation needs a PEB export plus an ECADSTAR PI re-simulation. Any later chapter
%      reusing the word "feasible" must carry the same caveat.
%  11. Known audit finding left UNFIXED (minor, out of evidence scope): the Stage-2 gate
%      for the diversity and occupancy-confidence terms is reported here as "with more than
%      one seed", which is what the facts block records. The audit reports the code also
%      requires the OPTIMIZE_BATCH_PER_K flag (default True). Verify against
%      pipelines/latent/optimize.py:69,414 and amend subsec:latent_space_opt.
%  12. Known audit finding left UNFIXED (bibliography, not this file): placeholder author
%      fields in refs.bib -- see the citation verification queue above.
%  13. The Bogatin section numbers cited in prose (13.12-13.17, 13.22) come from the
%      publisher's table of contents, not from the pages themselves. The physics written
%      around them (series RLC, SRF, parallel resonance, loop inductance) is standard, but
%      the section-to-claim mapping should be checked once against the book.
%  14. Academic-register pass: 69 file:line / artifact-path references were removed from the
%      body and preserved in ch2-impl-index.json. Numbers, equations, equation labels,
%      \fillin markers and all \cite keys were carried over unchanged; only starred
%      subsubsection headings were renamed (5 of them), per instruction.
% ---------------------------------------------------------------
